\documentclass[12pt, a4paper]{ctexart}

%==================== 导言区：宏包加载 ====================
\usepackage{amsmath, amssymb, amsthm} % 数学公式与符号
\usepackage{geometry}                 % 页面设置
\usepackage{xcolor}                   % 颜色支持
\usepackage{fancyhdr}                 % 页眉页脚
\usepackage{enumitem}                 % 列表定制
\usepackage{tcolorbox}                %以此美化文本框
\usepackage{titlesec}                 % 标题格式调整

%==================== 页面布局设置 ====================
\geometry{left=2.5cm, right=2.5cm, top=2.5cm, bottom=2.5cm}

%==================== 自定义视觉样式 ====================

% 1. 定义分隔线样式
\newcommand{\TopSeparator}{
    \par\noindent\rule{\textwidth}{1.5pt}\vspace{0.5em} % 粗实线 ━━━━
}
\newcommand{\AnalysisSeparator}{
    \par\noindent\rule{\textwidth}{0.5pt}\vspace{0.2em}\hrule\vspace{0.5em} % 双线效果 ════
}

% 2. 定义红色解析环境
%在这个环境内的所有文字和公式都会自动变成红色
\newenvironment{RedAnalysis}{
    \par
    \color{red} % 设置后续字体为红色
}{
    \par
    \color{black} % 结束时恢复黑色
}

% 3. 标题格式微调
\titleformat{\section}{\large\bfseries\heiti}{}{0em}{}[\vspace{-0.5em}]
\titleformat{\subsection}{\bfseries\kaishu}{}{0em}{}

\newcommand{\choicesFour}[4]{
    \begin{enumerate}[label=\Alph*., itemsep=0pt, parsep=0pt, topsep=5pt, labelsep=0.5em]
        \begin{minipage}{0.22\linewidth} \item #1 \end{minipage}
        \begin{minipage}{0.22\linewidth} \item #2 \end{minipage}
        \begin{minipage}{0.22\linewidth} \item #3 \end{minipage}
        \begin{minipage}{0.22\linewidth} \item #4 \end{minipage}
    \end{enumerate}
}
\newcommand{\choicesTwo}[4]{
    \begin{enumerate}[label=\Alph*., itemsep=0pt, parsep=0pt, topsep=5pt, labelsep=0.5em]
        \begin{minipage}{0.45\linewidth} \item #1 \end{minipage}
        \begin{minipage}{0.45\linewidth} \item #2 \end{minipage}
        \begin{minipage}{0.45\linewidth} \item #3 \end{minipage}
        \begin{minipage}{0.45\linewidth} \item #4 \end{minipage}
    \end{enumerate}
}
\newcommand{\choices}[4]{
    \begin{enumerate}[label=\Alph*., itemsep=0pt, parsep=0pt, topsep=5pt, labelsep=0.5em]
        \item #1
        \item #2
        \item #3
        \item #4
    \end{enumerate}
}

%==================== 文档开始 ====================
\begin{document}

% 标题信息
\title{\textbf{第五章 中值定理与导数应用（解析）}}
\author{数学习题讲义}
\date{\today}
\maketitle

% --- 题目 1 ---
\TopSeparator
\section*{题目1}

\textbf{【题目重现】} \\
设 $ f(x)=x\ln x $ , 则 $ f(x) $ \underline{\hspace{1cm}}。
\choicesTwo
{在 $ (0, \frac{1}{e}) $ 内单调减少}
{在 $ (\frac{1}{e}, +\infty) $ 内单调减少}
{在 $ (0, +\infty) $ 内单调减少}
{在 $ (0, +\infty) $ 内单调增加}

\AnalysisSeparator

\begin{RedAnalysis}

\subsection*{一、逐步解析}

\textbf{第一步：问题分析} \\
题目要求判断函数 $f(x) = x \ln x$ 的单调性。判断单调性的标准方法是求导数，根据导数的正负来确定单调区间。

\textbf{第二步：思路构建} \\
1. 求出函数 $f(x)$ 的定义域。
2. 求出 $f'(x)$。
3. 令 $f'(x) = 0$ 求出驻点。
4. 判断在驻点两侧导数的符号，从而确定单调区间。

\textbf{第三步：详细步骤} \\
1. \textbf{定义域}：由于含有 $\ln x$，故 $x > 0$。
2. \textbf{求导}：
   \[ f'(x) = (x \ln x)' = 1 \cdot \ln x + x \cdot \frac{1}{x} = \ln x + 1 \]
3. \textbf{求驻点}：令 $f'(x) = 0$，即 $\ln x + 1 = 0$，解得 $x = e^{-1} = \frac{1}{e}$。
4. \textbf{判断符号}：
   - 当 $0 < x < \frac{1}{e}$ 时，$\ln x < -1$，故 $f'(x) < 0$，函数单调减少。
   - 当 $x > \frac{1}{e}$ 时，$\ln x > -1$，故 $f'(x) > 0$，函数单调增加。
5. \textbf{对照选项}：
   A选项说在 $(0, \frac{1}{e})$ 内单调减少，符合计算结果。

\textbf{第四步：最终答案} \\
A

\subsection*{二、核心公式}
1. \textbf{乘积求导法则}：$(uv)' = u'v + uv'$
2. \textbf{单调性判别}：
   - $f'(x) > 0 \Rightarrow$ 单调增加
   - $f'(x) < 0 \Rightarrow$ 单调减少

\subsection*{三、考点分析}
\begin{itemize}
    \item \textbf{知识点归类}：导数的应用——利用导数判断函数的单调性。
    \item \textbf{易错点提示}：
       - 忽略定义域 $x>0$。
       - 对数函数的求导和不等式求解不熟练。
\end{itemize}

\end{RedAnalysis}

\vspace{2em}

% --- 题目 2 ---
\TopSeparator
\section*{题目2}

\textbf{【题目重现】} \\
曲线 $ y=6x-24x^{2}+x^{4} $ 的凸区间是 \underline{\hspace{1cm}}。
\choicesFour{$ (-2,2) $}{$ (-\infty,0) $}{$ (0,+\infty) $}{$ (-\infty,+\infty) $}

\AnalysisSeparator

\begin{RedAnalysis}

\subsection*{一、逐步解析}

\textbf{第一步：问题分析} \\
题目要求求曲线的凸区间。曲线的凹凸性由二阶导数的符号决定。
注意：教材关于“凹凸”的定义可能不同（有的教材称 $y'' > 0$ 为凹，有的称为凸）。
\textbf{通用标准（同济教材）}：
- $y'' > 0$：凹（向上弯曲）
- $y'' < 0$：凸（向下弯曲）
本题我们按通常标准：求凸区间即求 $y'' < 0$ 的区间。

\textbf{第二步：思路构建} \\
1. 求一阶导数 $y'$。
2. 求二阶导数 $y''$。
3. 解不等式 $y'' < 0$。

\textbf{第三步：详细步骤} \\
1. \textbf{求一阶导数}：
   \[ y' = 6 - 48x + 4x^3 \]
2. \textbf{求二阶导数}：
   \[ y'' = -48 + 12x^2 = 12(x^2 - 4) \]
3. \textbf{解不等式}：要求凸区间，即令 $y'' < 0$：
   \[ 12(x^2 - 4) < 0 \Rightarrow x^2 < 4 \Rightarrow -2 < x < 2 \]

\textbf{第四步：最终答案} \\
A

\subsection*{二、核心公式}
1. \textbf{凹凸性判别}：
   - $y'' > 0 \Rightarrow$ 凹曲线（开口向上）
   - $y'' < 0 \Rightarrow$ 凸曲线（开口向下）

\subsection*{三、考点分析}
\begin{itemize}
    \item \textbf{知识点归类}：导数的应用——曲线的凹凸性。
    \item \textbf{易错点提示}：凹凸性的定义容易混淆，需确认教材标准。通常考试中“凸”指向上凸起（即$y'' < 0$）。
\end{itemize}

\end{RedAnalysis}

\vspace{2em}

% --- 题目 3 ---
\TopSeparator
\section*{题目3}

\textbf{【题目重现】} \\
函数 $ y = \sin x $ 在区间 $ [0, \pi] $ 上满足罗尔定理的 $ \xi = $ \underline{\hspace{1cm}}。
\choicesFour{0}{$ \frac{\pi}{4} $}{$ \frac{\pi}{2} $}{$ \pi $}

\AnalysisSeparator

\begin{RedAnalysis}

\subsection*{一、逐步解析}

\textbf{第一步：问题分析} \\
罗尔定理要求函数在闭区间连续、开区间可导，且端点函数值相等。结论是存在 $\xi \in (a, b)$ 使得 $f'(\xi) = 0$。

\textbf{第二步：思路构建} \\
1. 验证条件：$f(0) = \sin 0 = 0$，$f(\pi) = \sin \pi = 0$，满足端点值相等。
2. 求导：$f'(x) = \cos x$。
3. 令导数为0：求 $\cos \xi = 0$ 在 $(0, \pi)$ 内的解。

\textbf{第三步：详细步骤} \\
1. $f'(x) = \cos x$。
2. 令 $f'(\xi) = 0$，即 $\cos \xi = 0$。
3. 在区间 $(0, \pi)$ 内，满足 $\cos \xi = 0$ 的值仅有 $\xi = \frac{\pi}{2}$。

\textbf{第四步：最终答案} \\
C

\subsection*{二、核心公式}
1. \textbf{罗尔定理}：若 $f(x) \in C[a,b]$，$f(x) \in D(a,b)$，且 $f(a)=f(b)$，则 $\exists \xi \in (a,b)$，使得 $f'(\xi)=0$。

\subsection*{三、考点分析}
\begin{itemize}
    \item \textbf{知识点归类}：中值定理——罗尔定理的验证与应用。
    \item \textbf{易错点提示}：$\xi$ 必须在开区间 $(0, \pi)$ 内部。
\end{itemize}

\end{RedAnalysis}

\vspace{2em}

% --- 题目 4 ---
\TopSeparator
\section*{题目4}

\textbf{【题目重现】} \\
曲线 $ y=\frac{x^{2}}{1+x} $ 的铅垂渐近线的方程是 \underline{\hspace{1cm}}。
\choicesFour{$ y = -1 $}{$ y = 1 $}{$ x = -1 $}{$ x = 1 $}

\AnalysisSeparator

\begin{RedAnalysis}

\subsection*{一、逐步解析}

\textbf{第一步：问题分析} \\
铅垂渐近线（垂直渐近线）通常出现在函数的间断点处。如果是分式函数，通常是分母为零的点。

\textbf{第二步：思路构建} \\
1. 找定义域的间断点。
2. 验证该点左右极限是否为无穷大。

\textbf{第三步：详细步骤} \\
1. 函数 $y=\frac{x^{2}}{1+x}$ 的定义域为 $x \neq -1$。可能的渐近线位置是 $x = -1$。
2. 计算极限：
   \[ \lim_{x \to -1} \frac{x^2}{1+x} = \lim_{x \to -1} \frac{1}{1+x} = \infty \]
   所以 $x = -1$ 是铅垂渐近线。

\textbf{第四步：最终答案} \\
C

\subsection*{二、核心公式}
1. \textbf{垂直渐近线判定}：若 $\lim_{x \to x_0} f(x) = \infty$，则 $x = x_0$ 是垂直渐近线。

\subsection*{三、考点分析}
\begin{itemize}
    \item \textbf{知识点归类}：导数的应用——函数的渐近线。
    \item \textbf{易错点提示}：垂直渐近线对应 $x=x_0$，水平渐近线对应 $y=A$。
\end{itemize}

\end{RedAnalysis}

\vspace{2em}

% --- 题目 5 ---
\TopSeparator
\section*{题目5}

\textbf{【题目重现】} \\
函数 $ f(x)=ax^{2}+b $ 在区间 $ (0,+\infty) $ 内单调增加，则 $a,b$ 应满足 \underline{\hspace{1cm}}。
\choicesTwo
{$a<0,b=0$}
{$a>0,b$ 为任意实数}
{$a<0,b\neq0$}
{$a<0,b$ 为任意实数}

\AnalysisSeparator

\begin{RedAnalysis}

\subsection*{一、逐步解析}

\textbf{第一步：问题分析} \\
已知函数在 $(0, +\infty)$ 单调增加，利用导数大于等于0来确定参数条件。

\textbf{第二步：思路构建} \\
1. 求导 $f'(x)$。
2. 令 $f'(x) \ge 0$ 在 $(0, +\infty)$ 恒成立（且不恒为0）。

\textbf{第三步：详细步骤} \\
1. 求导：$f'(x) = 2ax$。
2. 要求在 $(0, +\infty)$ 单调增加，即当 $x > 0$ 时，$f'(x) = 2ax > 0$。
3. 因为 $x > 0$，要使 $2ax > 0$，必须 $a > 0$。
4. 参数 $b$ 是常数项，求导后为0，不影响导数符号，所以 $b$ 可以是任意实数。

\textbf{第四步：最终答案} \\
B

\subsection*{二、核心公式}
1. \textbf{单调性与导数}：$f(x)$ 单增 $\Leftrightarrow f'(x) \ge 0$。

\subsection*{三、考点分析}
\begin{itemize}
    \item \textbf{知识点归类}：导数的应用——参数范围讨论。
\end{itemize}

\end{RedAnalysis}

\vspace{2em}

% --- 题目 6 ---
\TopSeparator
\section*{题目6}

\textbf{【题目重现】} \\
若在区间 $I$ 上， $ f'(x)>0, f''(x)<0 $ , 则曲线 $ y=f(x) $ 在 $I$ 是 \underline{\hspace{1cm}}。
\choicesTwo
{单调减少且为凹弧}
{单调减少且为凸弧}
{单调增加且为凹弧}
{单调增加且为凸弧}

\AnalysisSeparator

\begin{RedAnalysis}

\subsection*{一、逐步解析}

\textbf{第一步：问题分析} \\
直接根据导数符号的几何意义判断。

\textbf{第二步：思路构建} \\
- $f'(x) > 0$ 对应单调性。
- $f''(x) < 0$ 对应凹凸性。

\textbf{第三步：详细步骤} \\
1. $f'(x) > 0 \Rightarrow$ 函数单调增加。
2. $f''(x) < 0 \Rightarrow$ 曲线是凸的（向上凸起/开口向下）。

\textbf{第四步：最终答案} \\
D

\subsection*{二、核心公式}
1. \textbf{几何意义}：
   - $f' > 0$：增； $f' < 0$：减。
   - $f'' > 0$：凹； $f'' < 0$：凸。

\subsection*{三、考点分析}
\begin{itemize}
    \item \textbf{知识点归类}：导数应用——函数性态与导数符号的关系。
\end{itemize}

\end{RedAnalysis}

\vspace{2em}

% --- 题目 7 ---
\TopSeparator
\section*{题目7}

\textbf{【题目重现】} \\
关于曲线 $ y=\frac{2x^{3}}{(1-x)^{2}} $ ，以下说法正确的是 \underline{\hspace{1cm}}。
\choicesTwo
{既有水平渐近线，又有垂直渐近线}
{只有水平渐近线}
{有垂直渐近线 $ x = 1 $}
{没有渐近线}

\AnalysisSeparator

\begin{RedAnalysis}

\subsection*{一、逐步解析}

\textbf{第一步：问题分析} \\
分别检查垂直渐近线、水平渐近线和斜渐近线。

\textbf{第二步：思路构建} \\
1. **垂直渐近线**：找分母为0的点，检查极限。
2. **水平渐近线**：检查 $x \to \infty$ 时的极限。

\textbf{第三步：详细步骤} \\
1. **垂直渐近线**：
   定义域 $x \neq 1$。
   \[ \lim_{x \to 1} \frac{2x^3}{(1-x)^2} = \frac{2}{0^+} = +\infty \]
   所以 $x = 1$ 是垂直渐近线。

2. **水平渐近线**：
   \[ \lim_{x \to \infty} \frac{2x^3}{(1-x)^2} = \lim_{x \to \infty} \frac{2x^3}{x^2 - 2x + 1} = \lim_{x \to \infty} 2x = \infty \]
   极限不存在，所以没有水平渐近线。

3. **结论**：只有垂直渐近线 $x=1$（虽然可能有斜渐近线，但选项只讨论了水平和垂直，或者笼统的“有无”）。
   选项A错误（无水平）。
   选项B错误。
   选项C正确。
   选项D错误。

\textbf{第四步：最终答案} \\
C

\subsection*{二、核心公式}
1. \textbf{无穷远极限}：分子次数高于分母次数时，极限为无穷大，无水平渐近线。

\subsection*{三、考点分析}
\begin{itemize}
    \item \textbf{知识点归类}：导数应用——渐近线判定。
\end{itemize}

\end{RedAnalysis}

\vspace{2em}

% --- 题目 8 ---
\TopSeparator
\section*{题目8}

\textbf{【题目重现】} \\
在 $ (a,b) $ 区间内任意一点，函数 $ f(x) $ 的曲线弧总位于其切线的上方，则该曲线在 $ (a,b) $ 内是 \underline{\hspace{1cm}}。
\choicesFour{凹}{凸}{单调上升}{单调下降}

\AnalysisSeparator

\begin{RedAnalysis}

\subsection*{一、逐步解析}

\textbf{第一步：问题分析} \\
这是一道概念题，考查曲线凹凸性的几何特征。

\textbf{第二步：详细步骤} \\
如果曲线弧位于切线上方，说明曲线是“向下弯曲”或者是“开口向上”的形状。
在同济教材中：
- 曲线位于切线上方 $\Rightarrow$ 凹弧（$f''(x) > 0$）。
- 曲线位于切线下方 $\Rightarrow$ 凸弧（$f''(x) < 0$）。

\textbf{第四步：最终答案} \\
A

\subsection*{二、核心公式}
1. \textbf{凹凸性几何特征}：
   - 凹：切线上方。
   - 凸：切线下方。

\subsection*{三、考点分析}
\begin{itemize}
    \item \textbf{知识点归类}：导数应用——凹凸性的几何定义。
\end{itemize}

\end{RedAnalysis}

\vspace{2em}

% --- 题目 9 ---
\TopSeparator
\section*{题目9}

\textbf{【题目重现】} \\
当 $ x=\frac{\pi}{4} $ 时，函数 $ f(x)=a\cos x-\frac{1}{4}\cos 4x $ 取得极值，则 $a=$ \underline{\hspace{2cm}}。

\AnalysisSeparator

\begin{RedAnalysis}

\subsection*{一、逐步解析}

\textbf{第一步：问题分析} \\
函数在某点取得极值，如果该点可导，则该点处的导数必为0。

\textbf{第二步：思路构建} \\
1. 求 $f'(x)$。
2. 利用 $f'(\frac{\pi}{4}) = 0$ 建立方程解出 $a$。

\textbf{第三步：详细步骤} \\
1. \textbf{求导}：
   \[ f'(x) = -a\sin x - \frac{1}{4}(-\sin 4x) \cdot 4 = -a\sin x + \sin 4x \]
2. \textbf{代入条件}：
   因为在 $x = \frac{\pi}{4}$ 处取得极值，所以 $f'(\frac{\pi}{4}) = 0$。
   \[ -a\sin(\frac{\pi}{4}) + \sin(4 \cdot \frac{\pi}{4}) = 0 \]
   \[ -a \cdot \frac{\sqrt{2}}{2} + \sin(\pi) = 0 \]
   \[ -a \cdot \frac{\sqrt{2}}{2} + 0 = 0 \]
   \[ a = 0 \]

\textbf{第四步：最终答案} \\
0

\subsection*{二、核心公式}
1. \textbf{费马引理}：可导函数在极值点处的导数为0。

\subsection*{三、考点分析}
\begin{itemize}
    \item \textbf{知识点归类}：导数应用——极值存在的必要条件。
\end{itemize}

\end{RedAnalysis}

\vspace{2em}

% --- 题目 10 ---
\TopSeparator
\section*{题目10}

\textbf{【题目重现】} \\
函数 $ f(x)=\ln x $ 在区间 $[1,2]$ 上满足拉格朗日中值公式的 $ \xi= $ \underline{\hspace{2cm}}。

\AnalysisSeparator

\begin{RedAnalysis}

\subsection*{一、逐步解析}

\textbf{第一步：问题分析} \\
拉格朗日中值定理公式：$f(b) - f(a) = f'(\xi)(b-a)$。

\textbf{第二步：思路构建} \\
1. 计算 $f(b), f(a), b-a$。
2. 计算 $f'(\xi)$。
3. 解方程求 $\xi$。

\textbf{第三步：详细步骤} \\
1. 此处 $a=1, b=2$。
   $f(2) = \ln 2, f(1) = \ln 1 = 0$。
   $b-a = 2-1 = 1$。
   左边 $= \ln 2 - 0 = \ln 2$。
2. $f'(x) = \frac{1}{x} \Rightarrow f'(\xi) = \frac{1}{\xi}$。
   右边 $= \frac{1}{\xi} \cdot 1 = \frac{1}{\xi}$。
3. 令左边=右边：
   $\ln 2 = \frac{1}{\xi} \Rightarrow \xi = \frac{1}{\ln 2} = \log_2 e$。

\textbf{第四步：最终答案} \\
$\frac{1}{\ln 2}$ （或 $\log_2 e$）

\subsection*{二、核心公式}
1. \textbf{拉格朗日中值定理}：$f(b)-f(a) = f'(\xi)(b-a)$。

\subsection*{三、考点分析}
\begin{itemize}
    \item \textbf{知识点归类}：中值定理——拉格朗日中值定理的应用。
\end{itemize}

\end{RedAnalysis}

\vspace{2em}

% --- 题目 11 ---
\TopSeparator
\section*{题目11}

\textbf{【题目重现】} \\
设函数 $ f(x) $ 在 $ [1,2] $ 上连续，在 $ (1,2) $ 内可导，且 $ f(1) = 4f(2) $ 。
证明：存在一点 $ \xi \in (1,2) $ ，使得 $ 2f(\xi) + \xi f'(\xi) = 0 $ 。

\AnalysisSeparator

\begin{RedAnalysis}

\subsection*{一、逐步解析}

\textbf{第一步：问题分析} \\
证明含有 $f(\xi)$ 和 $f'(\xi)$ 的等式，通常需要构造辅助函数，然后使用罗尔定理。

\textbf{第二步：思路构建} \\
1. 观察目标式子：$2f(\xi) + \xi f'(\xi) = 0$。
2. 尝试还原原函数：$2f(x) + x f'(x)$ 看起来像 $(x^2 f(x))'$ 的一部分？
   $(x^2 f(x))' = 2x f(x) + x^2 f'(x) = x(2f(x) + xf'(x))$。
   这正是我们需要的结构（提取因子 $x \neq 0$ 后）。
   所以辅助函数设为 $F(x) = x^2 f(x)$。
3. 验证端点值：检查 $F(1)$ 和 $F(2)$ 是否相等。

\textbf{第三步：详细步骤} \\
1. \textbf{构造辅助函数}：
   令 $F(x) = x^2 f(x)$。
2. \textbf{验证罗尔定理条件}：
   - $F(x)$ 在 $[1,2]$ 上连续，在 $(1,2)$ 内可导（由 $f(x)$ 继承）。
   - 计算端点值：
     $F(1) = 1^2 \cdot f(1) = f(1)$
     $F(2) = 2^2 \cdot f(2) = 4f(2)$
   - 已知条件 $f(1) = 4f(2)$，所以 $F(1) = F(2)$。
3. \textbf{应用罗尔定理}：
   由罗尔定理，存在 $\xi \in (1,2)$ 使得 $F'(\xi) = 0$。
   \[ F'(x) = 2x f(x) + x^2 f'(x) \]
   所以 $2\xi f(\xi) + \xi^2 f'(\xi) = 0$。
4. \textbf{变形}：
   因为 $\xi \in (1,2)$，所以 $\xi \neq 0$。等式两边同时除以 $\xi$ 得：
   \[ 2f(\xi) + \xi f'(\xi) = 0 \]
   证毕。

\textbf{第四步：最终答案} \\
见上述证明过程。

\subsection*{二、核心公式}
1. \textbf{罗尔定理}：$F(a)=F(b) \Rightarrow \exists \xi, F'(\xi)=0$。
2. \textbf{辅助函数构造技巧}：观察导数形式反推原函数。
   - $f' + g' \to f+g$
   - $f'g + fg' \to fg$

\subsection*{三、考点分析}
\begin{itemize}
    \item \textbf{知识点归类}：中值定理——构造辅助函数证明等式。
    \item \textbf{技巧}：常见的辅助函数构造：
       - $f'(x) + k f(x) = 0 \Rightarrow F(x) = e^{kx}f(x)$
       - $xf'(x) + n f(x) = 0 \Rightarrow F(x) = x^n f(x)$
\end{itemize}

\end{RedAnalysis}

\vspace{2em}

% --- 题目 12 ---
\TopSeparator
\section*{题目12}

\textbf{【题目重现】} \\
用拉格朗日中值定理证明: 当 $x>0$ 时， $ \frac{x}{1+x}<\ln(1+x)<x $ 。

\AnalysisSeparator

\begin{RedAnalysis}

\subsection*{一、逐步解析}

\textbf{第一步：问题分析} \\
题目明确要求使用拉格朗日中值定理。不等式中间是 $\ln(1+x)$，可以看作函数 $f(t) = \ln(1+t)$ 在区间 $[0, x]$ 上的增量。

\textbf{第二步：思路构建} \\
1. 设 $f(t) = \ln(1+t)$。
2. 在区间 $[0, x]$ 上应用拉格朗日中值定理（其中 $x>0$）。
3. $f(x) - f(0) = f'(\xi)(x-0)$，其中 $0 < \xi < x$。
4. 利用 $\xi$ 的范围放大缩小来证明不等式。

\textbf{第三步：详细步骤} \\
1. 设 $f(t) = \ln(1+t)$。因为 $x>0$，在区间 $[0, x]$ 上 $f(t)$ 连续且可导。
2. 由拉格朗日中值定理：
   \[ \ln(1+x) - \ln(1+0) = \frac{1}{1+\xi} \cdot (x - 0), \quad 0 < \xi < x \]
   即 $\ln(1+x) = \frac{x}{1+\xi}$。
3. 因为 $0 < \xi < x$：
   - 分母：$1 < 1+\xi < 1+x$。
   - 倒数：$\frac{1}{1+x} < \frac{1}{1+\xi} < 1$。
   - 乘以 $x$ (因 $x>0$ 不等号不变)：$\frac{x}{1+x} < \frac{x}{1+\xi} < x$。
4. 代回 $\frac{x}{1+\xi} = \ln(1+x)$，得：
   \[ \frac{x}{1+x} < \ln(1+x) < x \]
   证毕。

\textbf{第四步：最终答案} \\
见上述证明过程。

\subsection*{二、核心公式}
1. \textbf{拉格朗日中值定理}：$f(b)-f(a) = f'(\xi)(b-a)$。

\subsection*{三、考点分析}
\begin{itemize}
    \item \textbf{知识点归类}：中值定理——利用中值定理证明不等式。
\end{itemize}

\end{RedAnalysis}

\vspace{2em}

% --- 题目 13 ---
\TopSeparator
\section*{题目13}

\textbf{【题目重现】} \\
证明：当 $x>0$ 时， $ \ln(1+x)>x-\frac{1}{2}x^{2} $ 。

\AnalysisSeparator

\begin{RedAnalysis}

\subsection*{一、逐步解析}

\textbf{第一步：问题分析} \\
这是一个单变量不等式证明。标准方法是构造辅助函数，利用导数判断单调性。

\textbf{第二步：思路构建} \\
1. 构造函数 $F(x) = \ln(1+x) - (x - \frac{1}{2}x^2)$。
2. 证明 $F(x)$ 在 $x>0$ 时单调递增。
3. 验证 $F(0) = 0$，从而推出 $F(x) > 0$。

\textbf{第三步：详细步骤} \\
1. \textbf{令} $F(x) = \ln(1+x) - x + \frac{1}{2}x^2$。
2. \textbf{求导}：
   \[ F'(x) = \frac{1}{1+x} - 1 + x = \frac{1 - (1+x) + x(1+x)}{1+x} = \frac{1-1-x+x+x^2}{1+x} = \frac{x^2}{1+x} \]
3. \textbf{判断符号}：
   因为 $x>0$，所以 $x^2 > 0, 1+x > 0$，故 $F'(x) > 0$。
4. \textbf{单调性结论}：即 $F(x)$ 在 $(0, +\infty)$ 上单调增加。
5. \textbf{比较初值}：$F(0) = \ln 1 - 0 + 0 = 0$。
   所以当 $x>0$ 时，$F(x) > F(0) = 0$。
   即 $\ln(1+x) - x + \frac{1}{2}x^2 > 0 \Rightarrow \ln(1+x) > x - \frac{1}{2}x^2$。

\textbf{第四步：最终答案} \\
见上述证明过程。

\subsection*{二、核心公式}
1. \textbf{不等式证明通用法}：构造 $F(x) = \text{左} - \text{右}$，证 $F'(x)>0$ 且 $F(0) \ge 0$。

\subsection*{三、考点分析}
\begin{itemize}
    \item \textbf{知识点归类}：导数应用——利用单调性证明不等式。
\end{itemize}

\end{RedAnalysis}

\vspace{2em}

% --- 题目 14 ---
\TopSeparator
\section*{题目14}

\textbf{【题目重现】} \\
求函数 $ y = x^{3} - 3x^{2} - 9x + 5 $ 的单调区间与极值。

\AnalysisSeparator

\begin{RedAnalysis}

\subsection*{一、逐步解析}

\textbf{第一步：问题分析} \\
题目要求两个内容：单调区间、极值。都需要通过一阶导数来求解。

\textbf{第二步：思路构建} \\
1. 求定义域（此处为 $\mathbb{R}$）。
2. 求 $y'$。
3. 令 $y'=0$ 求驻点。
4. 列表分析驻点两侧导数符号，得出单调区间和极值。

\textbf{第三步：详细步骤} \\
1. \textbf{求导}：$y' = 3x^2 - 6x - 9 = 3(x^2 - 2x - 3) = 3(x-3)(x+1)$。
2. \textbf{驻点}：令 $y'=0$，解得 $x_1 = -1, x_2 = 3$。
3. \textbf{列表分析}：（列表在此处用文字描述）
   - 当 $x < -1$ 时，$(x-3)<0, (x+1)<0$，则 $y' > 0$，单调增加。
   - 当 $-1 < x < 3$ 时，$(x-3)<0, (x+1)>0$，则 $y' < 0$，单调减少。
   - 当 $x > 3$ 时，$(x-3)>0, (x+1)>0$，则 $y' > 0$，单调增加。
4. \textbf{结论}：
   - **单调增区间**：$(-\infty, -1)$ 和 $(3, +\infty)$。
   - **单调减区间**：$(-1, 3)$。
5. \textbf{极值}：
   - $x = -1$ 处由增变减，极大值：
     $y(-1) = (-1)^3 - 3(-1)^2 - 9(-1) + 5 = -1 - 3 + 9 + 5 = 10$。
   - $x = 3$ 处由减变增，极小值：
     $y(3) = 3^3 - 3(3^2) - 9(3) + 5 = 27 - 27 - 27 + 5 = -22$。

\textbf{第四步：最终答案} \\
单调增区间：$(-\infty, -1) \cup (3, +\infty)$；单调减区间：$(-1, 3)$。
极大值 10；极小值 -22。

\subsection*{二、核心公式}
1. \textbf{极值第一充分条件}：导号由正变负取极大，由负变正取极小。

\subsection*{三、考点分析}
\begin{itemize}
    \item \textbf{知识点归类}：导数应用——单调性和极值的计算。
    \item \textbf{易错点提示}：计算函数值时容易算错，务必代回原函数检验。
\end{itemize}

\end{RedAnalysis}

\vspace{2em}

% --- 题目 15 ---
\TopSeparator
\section*{题目15}

\textbf{【题目重现】} \\
生产某种设备固定成本 1000 万元，每生产一台成本增加 20 万元，已知需求价格函数 $ P(Q) = 200 - Q $ ，问销量 $Q$ 为多少时，总利润 $L$ 最大，最大利润是多少？

\AnalysisSeparator

\begin{RedAnalysis}

\subsection*{一、逐步解析}

\textbf{第一步：问题分析} \\
这是一个经济应用题。
利润 = 收入 - 成本。
- 收入 $R(Q) = P(Q) \cdot Q$
- 成本 $C(Q) = \text{固定成本} + \text{变动成本}$
- 利润 $L(Q) = R(Q) - C(Q)$

\textbf{第二步：思路构建} \\
1. 写出 $C(Q)$ 和 $R(Q)$ 表达式。
2. 写出 $L(Q)$ 表达式。
3. 求 $L'(Q)$ 并令其为0，求驻点。
4. 验证是否为最大值点（实际问题通常只有一个驻点也就是最值点）。

\textbf{第三步：详细步骤} \\
1. \textbf{建立函数}：
   - 成本函数：$C(Q) = 1000 + 20Q$。
   - 收益函数：$R(Q) = P \cdot Q = (200 - Q)Q = 200Q - Q^2$。
   - 利润函数：$L(Q) = R(Q) - C(Q) = (200Q - Q^2) - (1000 + 20Q) = -Q^2 + 180Q - 1000$。
2. \textbf{求导求极值}：
   - $L'(Q) = -2Q + 180$。
   - 令 $L'(Q) = 0 \Rightarrow 2Q = 180 \Rightarrow Q = 90$。
3. \textbf{检验}：
   - $L''(Q) = -2 < 0$，说明曲线开口向下，驻点确实是最大值点。
4. \textbf{求最大利润}：
   - $L(90) = -(90)^2 + 180(90) - 1000 = -8100 + 16200 - 1000 = 8100 - 1000 = 7100$（万元）。

\textbf{第四步：最终答案} \\
当销量 $Q=90$ 时利润最大，最大利润为 7100 万元。

\subsection*{二、核心公式}
1. \textbf{经济学公式}：利润 = 收益 - 成本。
2. \textbf{最值应用}：求导找驻点。

\subsection*{三、考点分析}
\begin{itemize}
    \item \textbf{知识点归类}：导数应用——经济学中的最优化问题。
\end{itemize}

\end{RedAnalysis}

\end{document}
